\lecture{11}
\title{MLE Asymptotic Normality}

\begin{document}

\subsection{Defining Score}

We claim the MLE also achieves asymptotic normality.  In fact,
we claim we even know its asymptotic variance.
It turns out, however, that we'll need a little more setup before we get to asymptotic variance.

\textbf{Definition.} Consider a \underline{statistical model} with a parameter $\theta$
and varying PDFs given by $f_{\theta}(x)$, dependent on $\theta$.
Given an observation $x$ and a proposed value $\theta$ of the parameter,
the \textbf{score} is:
\[ s(x; \theta) := \frac{\partial}{\partial\theta} \log f_{\theta}(x)
= \frac{\frac{\partial}{\partial\theta} f_{\theta}(x)}{f_{\theta}(x)}. \]
The RHS comes from a single application of the chain rule.
When $s(x; \theta)$ is positive, it means that
$x$ would be better explained if the proposed $\theta$ were increased, and vice versa.

\begin{remark}
    The notation is slightly confusing;
    as a simpler example, consider defining $g(x) := \frac{d}{dx} x^2$ and concluding $g(5) = 10$.
\end{remark}

\begin{theorem}{Zero Mean Score}
    We have $\mathbb{E}_{\theta^*}[s(X; \theta^*)] = 0$,
    where the $\mathbb{E}_{\theta^*}$ notation indicates that
    $X$ is sampled from the distribution $\mathbb{P}_{\theta^*}$;
    in other words, the proposed value $\theta^*$ is correct.

    In this case, the score has some positives and negatives, but expectedly is overall neutral
    in how the proposal $\theta^*$ should change.
\end{theorem}

\begin{proof}
    Just some integral manipulation.
    \[ \mathbb{E}_{\theta^*}[s(X; \theta^*)]
    = \int_{\mathcal{R}} f_{\theta^*}(x) \cdot
    \left[\frac{\frac{\partial}{\partial\theta} f_{\theta}(x)}{f_{\theta}(x)}\right]_{\theta = \theta^*}
    \mathrm{d}x
    = \int_{\mathcal{R}} \left[\frac{\partial}{\partial\theta} f_{\theta}(x)\right]_{\theta = \theta^*}
    \mathrm{d}x
    = \frac{\partial}{\partial\theta}
    \left[\int_{\mathcal{R}} f_{\theta}(x) \, \mathrm{d}x\right]_{\theta = \theta^*}
    = \frac{\partial}{\partial\theta} [1] = 0. ~~~ \blacksquare \]
\end{proof}

\subsection{Defining Fisher Information}

Note that the definition of score doesn't care whether $x$ is sampled
from a distribution that uses the proposed value $\theta$ for its parameter.
The definition of Fisher information, on the other hand, \emph{does} care.

\textbf{Definition.} Consider a \underline{statistical model} with a parameter $\theta$
and varying PDFs given by $f_{\theta}(x)$, dependent on $\theta$.
As a function of the true value of $\theta$, the \textbf{Fisher information}
of the resulting distribution is:
\[ I(\theta) := \mathbb{V}_{\theta}\left[s(X; \theta)\right]
= \mathbb{E}_{\theta}[s(X; \theta)^2]. \]
As usual, the $\mathbb{E}_{\theta}$ and $\mathbb{V}_{\theta}$ notation
indicates that $X \sim \mathbb{P}_{\theta}$; in other words, that $\theta$ is the ``true value''
of the parameter.  It's because of this and the previous theorem
that the above two expressions for $I(\theta)$ are equivalent.

When $I(\theta)$ is large,
it means a single observation $x$ will tend to suggest
very \emph{large} changes to the proposed $\theta$.
In other words, singular observations $x$ carry more information.

\begin{problem}
    Given the model $\{\mathcal{N}(\theta, 9) \mid \theta \in \mathbb{R}\}$,
    determine $I(\theta)$.  Do the same for $\{\mathcal{N}(\theta, 900) \mid \theta \in \mathbb{R}\}$.
\end{problem}

\begin{solution}
    Recall that the PDF of $X \sim \mathcal{N}(\theta, 9)$ (with $\sigma = 3$) is given by:
    \[ f_{\theta}(x) = \frac{1}{3\sqrt{2\pi}} \exp\left(-\frac{(x - \theta)^2}{2 \times 3^2}\right)
    ~ \implies ~ s(x; \theta) = \frac{\partial}{\partial\theta} \log f_{\theta}(x)
    = \frac{x - \theta}{3^2}. \]
    The Fisher information is therefore:
    \[ I(\theta) = \mathbb{V}_{\theta}\left[ \frac{X - \theta}{3^2} \right]
    = \left(\frac{1}{3^2}\right)^2 \mathbb{V}_{\theta}[X - \theta]
    = \left(\frac{1}{3^2}\right)^2(3^2) = \frac{1}{9}. \]
    It's a constant!
    Expectedly, the Fisher information for $\{\mathcal{N}(\theta, 900) \mid \theta \in \mathbb{R}\}$
    is constantly $\frac{1}{900}$.
    
    Notably, the Fisher information of the latter is smaller than the former.
    This should make sense; a single observation from
    $\{\mathcal{N}(\theta, 900) \mid \theta \in \mathbb{R}\}$
    tells you a lot less about $\theta$ than a single observation from
    $\{\mathcal{N}(\theta, 9) \mid \theta \in \mathbb{R}\}$ does.
\end{solution}

The Fisher information has one more alternative representation.

\begin{theorem}{Equivalent Fisher}
    Fisher information equals the negative expected curvature of the log likelihood; that is,
    \[ I(\theta) = -\mathbb{E}_{\theta}\left[
        \frac{\partial^2}{\partial\theta^2} \log f_{\theta}(X)
    \right]. \]
\end{theorem}

\begin{proof}
    Just a lot of calculus.  By the quotient and chain rules, we have:
    \[ \frac{\partial^2}{\partial\theta^2} \log f_{\theta}(X)
    = \frac{\partial}{\partial\theta} \frac{\frac{\partial}{\partial\theta} f_{\theta}(X)}{f_{\theta}(X)} 
    = \frac{f_{\theta}(X) \frac{\partial^2}{\partial\theta^2} f_{\theta}(X)
    - \left(\frac{\partial}{\partial\theta} f_{\theta}(X)\right)^2}{
        \left(f_{\theta}(X)\right)^2
    }
    = \frac{\frac{\partial^2}{\partial\theta^2} f_{\theta}(X)}{f_{\theta}(X)}
    - s(X; \theta)^2. \]
    The expectation of the second term of the RHS is $-I(\theta)$,
    so it suffices to show the expectation of the first term is zero.
    \[ \mathbb{E}_{\theta}\left[\frac{\frac{\partial^2}{\partial\theta^2} f_{\theta}(X)}{f_{\theta}(X)}\right]
    = \int_{\mathcal{R}} f_{\theta}(x) \cdot \left[
        \frac{\frac{\partial^2}{\partial\theta^2} f_{\theta}(x)}{f_{\theta}(x)}
    \right] \, \mathrm{d}x
    = \int_{\mathcal{R}} \left[\frac{\partial^2}{\partial\theta^2} f_{\theta}(x)\right] \, \mathrm{d}x
    = \frac{\partial^2}{\partial\theta^2} \left[\int_{\mathcal{R}} f_{\theta}(x) \, \mathrm{d}x\right]
    = \frac{\partial^2}{\partial\theta^2} [1] = 0. \]
    A very similar argument as before, actually\dots ~ $\blacksquare$
\end{proof}

Finally, when we have multiple samples,
we can talk about ``total Fisher information''.

\textbf{Definition.} Given $n$ samples $\{X_i\}_{i = 1}^n$,
recall $\ell_n(\theta) := \sum_{i = 1}^n \log f_{\theta}(X_i)$
is the log likelihood.
Then the \textbf{total Fisher information} gained from these samples is:
\[ I_n(\theta) :=
-\mathbb{E}_{\theta}\left[\frac{\partial^2}{\partial\theta^2}\ell_n(\theta)\right]
= n I(\theta). \]
These two expressions are equal because all the $X_i$ are sampled from
the same distribution $\mathbb{P}_{\theta}$, so applying linearity of expectation
yields $n$ copies of $I(\theta)$ summed together.

\subsection{MLE Asymptotic Normality}

Now we can finally get to asymptotic normality of the MLE.

\begin{theorem}{MLE Asymptotic Normality}
    Under mild regularity conditions,
    the MLE $\hat{\theta}_n$ has asymptotic variance $\frac{1}{I(\theta^*)}$; that is,
    \[ \sqrt{n}(\hat{\theta}_n - \theta^*) \rightsquigarrow
    \mathcal{N}\left(0, \frac{1}{I(\theta^*)}\right). \]
\end{theorem}

\begin{proof}
    Recall the log likelihood
    $\ell_n(\theta) := \sum_{i = 1}^n \log f_{\theta}(X_i)$.
    The definition of the MLE $\hat{\theta}_n$ is that it maximizes $\ell_n(\theta)$, meaning:
    \[ \frac{\partial}{\partial\theta} \ell_n(\hat{\theta}_n)
    = \sum_{i = 1}^n s(X_i; \hat{\theta}_n) = 0. \]
    Recall that we proved the consistency of $\hat{\theta}_n$ earlier,
    meaning $\hat{\theta}_n \convprob \theta^*$.
    Thus, it makes sense to take the Taylor approximation:
    \[ 0 = \sum_{i = 1}^n s(X_i; \hat{\theta}_n)
    \approx \sum_{i = 1}^n s(X_i; \theta^*) +
    (\hat{\theta}_n - \theta^*) \left[
        \frac{\partial}{\partial\theta} \sum_{i = 1}^n s(X_i; \theta)
    \right]_{\theta = \theta^*}. \]
    Moving terms around yields the approximation:
    \[ \sqrt{n}(\hat{\theta}_n - \theta^*)
    \approx - \left(\frac{\frac{1}{\sqrt{n}}\cdot \sum_{i = 1}^n s(X_i; \theta^*)}{
        \frac{1}{n} \cdot \sum_{i = 1}^n \left[\frac{\partial^2}{\partial\theta^2} \log f_{\theta}(X_i)
        \right]_{\theta = \theta^*}
    }\right). \]
    All that setup pays off here;
    it turns out the numerator and denominator of this fraction
    can be expressed very nicely.
    \begin{itemize}
        \item By ``Zero Mean Score'', the definition of Fisher information, and CLT,
        the numerator converges ($\rightsquigarrow$) to $\mathcal{N}(0, I(\theta^*))$.
        \item By ``Equivalent Fisher'' and LLN,
        the denominator converges ($\convprob$) to $-I(\theta^*)$.
    \end{itemize}
    Thus, by Slutsky's lemma, we have:
    \[ \sqrt{n}(\hat{\theta}_n - \theta^*) \rightsquigarrow
    -\left(\frac{\mathcal{N}(0, I(\theta^*))}{-I(\theta^*)}\right)
    = \mathcal{N}\left(0, \frac{1}{I(\theta^*)}\right). \]
    And that's it. ~ $\blacksquare$
\end{proof}

Knowing that the asymptotic variance is $\frac{1}{I(\theta^*)}$
is useful for constructing confidence intervals using the MLE $\hat{\theta}_n$,
especially when the variance of $\hat{\theta}$ would not otherwise be easy to compute directly.

\begin{remark}
    For the model $\{\mathrm{Ber}(p) \mid p \in (0, 1)\}$,
    the Fisher information is $I(p) = \frac{1}{p(1 - p)}$.
    CYU: Why is this expected?
\end{remark}

\subsection{The Cramér-Rao Lower Bound}

As the following theorem explains,
the MLE is in some sense the \emph{best estimator possible},
as long as it is unbiased.

\begin{theorem}{Cramér-Rao Lower Bound}
    Let $\hat{\theta}_n$ be any unbiased estimator of $\theta$.
    Then $\mathbb{V}_{\theta^*}[\hat{\theta}_n] \geq \frac{1}{nI(\theta^*)}$.
\end{theorem}

\begin{proof}
    The finish of the proof will use the following theorem.
    \begin{quote}
        \textbf{Theorem (Cauchy-Schwarz).}
        For any random variables $U$ and $V$,
        we have $\Cov(U, V)^2 \leq \Var(U) \cdot \Var(V)$.

        \emph{Proof:} The point is that $\Var(U - \lambda V) \geq 0$
        for all $\lambda$.  The LHS expands to a quadratic in $\lambda$,
        so the discriminant of this quadratic must be nonpositive.
        The nonpositivity of the discriminant simplifies to the above theorem.
        ~ $\square$
    \end{quote}
    We'll invoke the above theorem on
    $U := \hat{\theta}_n$ and $V := \sum_{i = 1}^n s(X_i; \theta^*)$.
    \begin{quote}
        \textbf{Lemma 1.} We have $\Var_{\theta^*}(V) = nI(\theta^*)$.

        \emph{Proof:} The $X_i$ are independent, so variances add,
        and then this follows straight from the definition of Fisher information. ~ $\square$
    \end{quote}
    Now for the hard part.
    \begin{quote}
        \textbf{Lemma 2.} We have $\Cov_{\theta^*}(U, V) = 1$.

        \emph{Proof:} Note that $\mathbb{E}_{\theta^*}\left[
            \sum_{i = 1}^n s(X_i; \theta^*)
        \right] = 0$ by ``Zero Mean Score''.
        Also note that we're considering
        the joint random variable $(X_1, \dots, X_n)$,
        so our integrals will need to use the joint PDFs $\prod_{i = 1}^n f_{\theta^*}(x_i)$.
        That said, the rest is computation.
        \begin{align*}
        \Cov_{\theta^*}(U, V)
        & = \mathbb{E}_{\theta^*}\left[\hat{\theta} \cdot \sum_{i = 1}^n s(X_i; \theta^*)\right] -
        \mathbb{E}_{\theta^*}[\hat{\theta}] \cdot
        \mathbb{E}_{\theta^*}\left[\sum_{i = 1}^n s(X_i; \theta^*)\right] \\
        & = \mathbb{E}_{\theta^*}\left[\hat{\theta} \cdot \sum_{i = 1}^n s(X_i; \theta^*)\right] \\
        & = \int \cdots \int_{\mathcal{R}} \left(\prod_{i = 1}^n f_{\theta^*}(x_i)\right) \cdot
        \hat{\theta}(x_1, x_2, \dots, x_n) \cdot \sum_{i = 1}^n s(x_i; \theta^*)
        \, \mathrm{d}x_1 \, \mathrm{d}x_2 \, \cdots \, \mathrm{d}x_n \\
        & = \int \cdots \int_{\mathcal{R}} \left(\prod_{i = 1}^n f_{\theta^*}(x_i)\right) \cdot
        \hat{\theta}(x_1, x_2, \dots, x_n) \cdot
        \frac{
            \left[\frac{\partial}{\partial\theta}
            \prod_{i = 1}^n f_{\theta}(x_i)\right]_{\theta = \theta^*}
        }{\prod_{i = 1}^n f_{\theta^*}(x_i)}
        \, \mathrm{d}x_1 \, \mathrm{d}x_2 \, \cdots \, \mathrm{d}x_n \\
        & = \frac{\partial}{\partial\theta} \left[
        \int \cdots \int_{\mathcal{R}} \left(\prod_{i = 1}^n f_{\theta}(x_i)\right) \cdot
        \hat{\theta}(x_1, x_2, \dots, x_n)
        \, \mathrm{d}x_1 \, \mathrm{d}x_2 \, \cdots \, \mathrm{d}x_n
        \right]_{\theta = \theta^*}.
        \end{align*}
        Here's the finish: $\hat{\theta}$ is \emph{unbiased}, so
        the integral above evaluates to just $\theta$,
        meaning the covariance is just
        $\frac{\partial}{\partial\theta}[\theta]_{\theta = \theta^*} = 1$.
        ~ $\square$
    \end{quote}
    Now just plug Lemma 1 and Lemma 2 into the Cauchy-Schwarz theorem and win. ~ $\blacksquare$
\end{proof}

\end{document}